\documentclass[12pt, a4paper]{article}

% ─── Codificación y lengua ───────────────────────────────────────────
\usepackage[utf8]{inputenc}
\usepackage[T1]{fontenc}
\usepackage[spanish, mexico]{babel}

% ─── Tipografía ──────────────────────────────────────────────────────
\usepackage{lmodern}
\usepackage{microtype}

% ─── Matemáticas ─────────────────────────────────────────────────────
\usepackage{amsmath, amssymb, amsthm}
\usepackage{mathtools}
\usepackage{bm}

% ─── Gráficos y colores ──────────────────────────────────────────────
\usepackage{graphicx}
\usepackage[table, xcdraw, dvipsnames]{xcolor}
\usepackage{tikz}
\usetikzlibrary{shapes.geometric, arrows.meta, positioning, fit, backgrounds, calc}

% ─── Tablas ──────────────────────────────────────────────────────────
\usepackage{booktabs}
\usepackage{tabularx}
\usepackage{multirow}
\usepackage{array}

% ─── Código fuente ───────────────────────────────────────────────────
\usepackage{listings}
\usepackage{lstautogobble}

% ─── Diseño de página ────────────────────────────────────────────────
\usepackage[top=2.5cm, bottom=2.5cm, left=3cm, right=2.5cm]{geometry}
\usepackage{setspace}
\onehalfspacing
\usepackage{parskip}
\setlength{\parindent}{0pt}
\setlength{\parskip}{6pt}

% ─── Encabezados y pie de página ─────────────────────────────────────
\usepackage{fancyhdr}
\pagestyle{fancy}
\fancyhf{}
\fancyhead[L]{\small\textit{UrbanFlow CDMX — Predicción Estocástica de Tiempos de Viaje}}
\fancyhead[R]{\small\thepage}
\renewcommand{\headrulewidth}{0.4pt}

% ─── Hiperenlaces ────────────────────────────────────────────────────
\usepackage[colorlinks=true, linkcolor=NavyBlue, citecolor=NavyBlue,
            urlcolor=NavyBlue, bookmarks=true]{hyperref}

% ─── Cajas y entornos especiales ─────────────────────────────────────
\usepackage{tcolorbox}
\tcbuselibrary{skins, breakable, theorems}

% ─── Títulos de sección ──────────────────────────────────────────────
\usepackage{titlesec}
\titleformat{\section}{\large\bfseries\color{NavyBlue}}{{\thesection}}{0.8em}{}[\color{NavyBlue}\titlerule]
\titleformat{\subsection}{\normalsize\bfseries\color{MidnightBlue}}{\thesubsection}{0.8em}{}
\titleformat{\subsubsection}{\normalsize\itshape\color{MidnightBlue}}{\thesubsubsection}{0.8em}{}

% ─── Algoritmos ──────────────────────────────────────────────────────
\usepackage[ruled, vlined, linesnumbered]{algorithm2e}
\SetAlgorithmName{Algoritmo}{algoritmo}{Índice de algoritmos}

% ─── Definiciones de teoremas ────────────────────────────────────────
\theoremstyle{definition}
\newtheorem{definicion}{Definición}[section]
\newtheorem{propiedad}{Propiedad}[section]
\newtheorem{observacion}{Observación}[section]

% ─── Configuración de listings (Python) ──────────────────────────────
\definecolor{codebg}{RGB}{248, 248, 252}
\definecolor{codecomment}{RGB}{100, 100, 150}
\definecolor{codestring}{RGB}{180, 40, 40}
\definecolor{codekeyword}{RGB}{0, 80, 160}

\lstdefinestyle{python}{
    language=Python,
    backgroundcolor=\color{codebg},
    basicstyle=\ttfamily\footnotesize,
    keywordstyle=\bfseries\color{codekeyword},
    stringstyle=\color{codestring},
    commentstyle=\itshape\color{codecomment},
    numberstyle=\tiny\color{gray},
    numbers=left,
    numbersep=8pt,
    stepnumber=1,
    breaklines=true,
    frame=single,
    framesep=4pt,
    rulecolor=\color{gray!40},
    showstringspaces=false,
    tabsize=4,
    autogobble=true,
    aboveskip=8pt,
    belowskip=8pt,
}
\lstset{style=python}

% ─── Colores personalizados ──────────────────────────────────────────
\definecolor{fluido}{RGB}{39, 174, 96}
\definecolor{lento}{RGB}{230, 126, 34}
\definecolor{congestionado}{RGB}{192, 57, 43}
\definecolor{headerblue}{RGB}{31, 60, 122}
\definecolor{lightblue}{RGB}{235, 242, 255}

% ─── Entornos tcolorbox personalizados ───────────────────────────────
\newtcolorbox{cajadefinicion}[1]{
    colback=lightblue, colframe=NavyBlue,
    fonttitle=\bfseries, title=#1,
    breakable, arc=4pt
}
\newtcolorbox{cajaresultado}[1]{
    colback=green!5, colframe=fluido!80!black,
    fonttitle=\bfseries, title=#1,
    breakable, arc=4pt
}

% ─── Metadatos del documento ─────────────────────────────────────────
\title{%
    {\LARGE \textbf{UrbanFlow CDMX}}\\[4pt]
    {\large Sistema de Predicción Estocástica de Tiempos de Viaje\\
    en la Zona Metropolitana del Valle de México}\\[6pt]
    {\normalsize mediante Cadenas de Markov y Simulación Monte Carlo}
}
\author{%
    \textbf{Carlos}\\[4pt]
    \small Diplomado en Ciencia de Datos\\
    \small Proyecto Integrador Final
}
\date{\small Marzo 2026}

% ════════════════════════════════════════════════════════════════════
\begin{document}
% ════════════════════════════════════════════════════════════════════

\maketitle
\thispagestyle{empty}

\begin{abstract}
\noindent
UrbanFlow CDMX es un sistema de predicción probabilística de tiempos de viaje
para la Zona Metropolitana del Valle de México (ZMVM) que combina
\textbf{cadenas de Markov de tiempo discreto} con \textbf{simulación Monte Carlo}
de $N = 10{,}000$ trayectorias por consulta. El sistema ingiere datos en tiempo
real de TomTom Traffic Stats API, TomTom Routing API y OpenWeatherMap, construye
un modelo estocástico del comportamiento del tráfico y entrega al usuario las
bandas de incertidumbre $P_{10}$, $P_{50}$ y $P_{90}$ del tiempo de viaje.
La métrica objetivo es un \textbf{MAPE $< 15\%$} para horizontes de predicción
de 30 minutos. A diferencia de los sistemas de navegación convencionales que
producen una estimación puntual, UrbanFlow CDMX cuantifica explícitamente la
incertidumbre inherente al tráfico urbano, permitiendo tomar decisiones racionales
de planeación de viaje.
\end{abstract}

\vspace{0.5cm}
\tableofcontents
\newpage

% ════════════════════════════════════════════════════════════════════
\section{Motivación y Contexto}
% ════════════════════════════════════════════════════════════════════

\subsection{El problema del tráfico en la ZMVM}

La Ciudad de México y su zona conurbada constituyen una de las áreas urbanas
más congestionadas del mundo. Según el \textit{TomTom Traffic Index 2024},
los conductores de la CDMX pierden en promedio \textbf{119 horas al año} en
congestionamientos, lo que la posiciona entre las diez ciudades con mayor
pérdida de tiempo por tráfico a nivel global. Este fenómeno tiene consecuencias
directas en productividad económica, emisiones de $\text{CO}_2$ y calidad de
vida de los más de 21 millones de personas que habitan la ZMVM.

\subsection{La limitación del enfoque determinístico}

Los sistemas de navegación convencionales (Google Maps, Waze) producen predicciones
de tiempo de viaje \textbf{puntuales}: un único número. Sin embargo, el tráfico
urbano es inherentemente \textbf{estocástico}: el mismo recorrido a la misma
hora puede variar 20–40 minutos dependiendo de incidentes viales, condiciones
climáticas y el comportamiento colectivo de los demás conductores.

\subsection{Propuesta de valor}

UrbanFlow CDMX reemplaza la predicción puntual por una
\textbf{distribución de probabilidad completa} del tiempo de viaje,
expresada mediante tres percentiles:

\begin{table}[h!]
\centering
\renewcommand{\arraystretch}{1.3}
\begin{tabular}{clp{8cm}}
\toprule
\textbf{Percentil} & \textbf{Escenario} & \textbf{Significado} \\
\midrule
$P_{10}$ & Optimista  & El 10\% de los días se llega antes de este tiempo \\
$P_{50}$ & Probable   & Mediana: tan probable llegar antes como después \\
$P_{90}$ & Pesimista  & El 90\% de los días se llega antes de este tiempo \\
\bottomrule
\end{tabular}
\caption{Bandas de incertidumbre del tiempo de viaje en UrbanFlow CDMX.}
\end{table}

% ════════════════════════════════════════════════════════════════════
\section{Arquitectura General del Sistema}
% ════════════════════════════════════════════════════════════════════

\subsection{Diagrama de componentes}

\begin{figure}[h!]
\centering
\begin{tikzpicture}[
    node distance=1.2cm and 0.8cm,
    every node/.style={font=\small},
    fuente/.style={rectangle, rounded corners=4pt, draw=NavyBlue, fill=lightblue,
                   minimum width=3cm, minimum height=1cm, text centered,
                   text width=2.8cm},
    modulo/.style={rectangle, rounded corners=4pt, draw=MidnightBlue,
                   fill=MidnightBlue!10, minimum width=3.5cm, minimum height=1cm,
                   text centered, text width=3.3cm},
    motor/.style={rectangle, rounded corners=4pt, draw=fluido!80!black,
                  fill=fluido!10, minimum width=3.5cm, minimum height=1cm,
                  text centered, text width=3.3cm},
    output/.style={rectangle, rounded corners=4pt, draw=congestionado!80!black,
                   fill=congestionado!10, minimum width=5cm, minimum height=1cm,
                   text centered, text width=4.8cm},
    flecha/.style={-{Stealth[length=6pt]}, thick, color=NavyBlue},
]

% Fuentes de datos
\node[fuente] (tomtom_t) {TomTom\\Traffic Stats};
\node[fuente, right=0.6cm of tomtom_t] (tomtom_r) {TomTom\\Routing API};
\node[fuente, right=0.6cm of tomtom_r] (owm) {OpenWeather\\Map API};

% Pipeline integrador
\node[modulo, below=1.2cm of tomtom_r] (pipeline)
    {\textbf{Pipeline Integrador}\\velocidades + clima\\$\rightarrow$ ContextoViaje};

% Motor estocástico
\node[motor, left=1.2cm of pipeline, yshift=-2.2cm] (markov)
    {\textbf{MarkovTrafficChain}\\$\hat{P}_{ij}$ estimada\\desde datos históricos};
\node[motor, right=1.2cm of pipeline, yshift=-2.2cm] (mc)
    {\textbf{MonteCarloEngine}\\$N=10{,}000$ trayectorias\\$P_{10}/P_{50}/P_{90}$};

% Output
\node[output, below=1.2cm of mc] (resultado)
    {\textbf{ResultadoSimulacion}\\$P_{10},\; P_{50},\; P_{90}$,\; $\bar{t}$,\; $\sigma_t$};
\node[output, below=1.2cm of markov] (dashboard)
    {\textbf{Dashboard Streamlit}\\Mapa · Gauge · Histograma\\Matriz de transición};

% Flechas
\draw[flecha] (tomtom_t) -- (pipeline);
\draw[flecha] (tomtom_r) -- (pipeline);
\draw[flecha] (owm) -- (pipeline);
\draw[flecha] (pipeline) -- (markov);
\draw[flecha] (pipeline) -- (mc);
\draw[flecha] (markov) -- (mc);
\draw[flecha] (mc) -- (resultado);
\draw[flecha] (resultado) -- (dashboard);
\draw[flecha] (markov) -- (dashboard);
\end{tikzpicture}
\caption{Arquitectura de componentes de UrbanFlow CDMX.}
\end{figure}

\subsection{Stack tecnológico}

\begin{table}[h!]
\centering
\renewcommand{\arraystretch}{1.3}
\begin{tabularx}{\textwidth}{llX}
\toprule
\textbf{Capa} & \textbf{Herramienta} & \textbf{Rol} \\
\midrule
Cómputo numérico & NumPy & Álgebra lineal vectorizada (Markov, Monte Carlo) \\
Datos tabulares  & Pandas & Manipulación de DataFrames de tráfico y clima \\
Geoespacial      & GeoPandas / Folium & Proyección EPSG:6372, mapa interactivo \\
ML / Forecasting & XGBoost, Prophet & Modelos complementarios de predicción \\
Frontend         & Streamlit & Dashboard interactivo en tiempo real \\
APIs externas    & TomTom, OWM & Velocidades en tiempo real y clima \\
Serialización    & \texttt{dataclasses}, JSON & Contratos tipados entre módulos \\
\bottomrule
\end{tabularx}
\caption{Stack tecnológico del sistema UrbanFlow CDMX.}
\end{table}

% ════════════════════════════════════════════════════════════════════
\section{Fuentes de Datos e Ingesta}
% ════════════════════════════════════════════════════════════════════

\subsection{TomTom Traffic Stats API}

Proporciona velocidades en tiempo real por coordenada. El cliente
\texttt{tomtom\_client.py} obtiene para cada punto del corredor:
$v_\text{actual}$ (km/h), $v_\text{libre}$ (km/h) y el
\textbf{ratio de congestión}:
\begin{equation}
    r = \frac{v_\text{actual}}{v_\text{libre}} \in (0, 1]
    \label{eq:ratio}
\end{equation}

El ratio $r$ es la señal principal para inferir el estado inicial de tráfico
que se pasa a la cadena de Markov (véase Sección~\ref{sec:pipeline}).

\subsection{TomTom Routing API}

Calcula la \textbf{distancia real por carretera} entre dos puntos, reemplazando
el estimador Haversine (distancia en línea recta). El endpoint utilizado es:
\begin{center}
\texttt{GET /routing/1/calculateRoute/\{lat1,lon1\}:\{lat2,lon2\}/json}
\end{center}

La respuesta incluye \texttt{lengthInMeters} y la geometría de la ruta
(polilínea de hasta 120 puntos lat/lon submuestreados) para renderizar en
el mapa Folium.

\subsubsection*{Fallback empírico de tortuosidad urbana}

Cuando la API no está disponible se aplica el factor de tortuosidad:
\begin{equation}
    d_\text{carretera} \approx d_\text{haversine} \times \kappa, \quad \kappa = 1.4
    \label{eq:tortuosidad}
\end{equation}
El factor $\kappa = 1.4$ está calibrado empíricamente para la ZMVM: la red vial
urbana tiene una tortuosidad media del 40\% sobre la distancia euclidiana.

\subsection{OpenWeatherMap API}

Proporciona temperatura, precipitación, velocidad del viento y condición
climática. El módulo \texttt{weather\_client.py} convierte la condición
meteorológica en un \textbf{factor climático multiplicativo} $\phi \in (0,1]$:

\begin{table}[h!]
\centering
\renewcommand{\arraystretch}{1.3}
\begin{tabular}{lcc}
\toprule
\textbf{Condición} & \textbf{Factor $\phi$} & \textbf{Reducción velocidad} \\
\midrule
Despejado    & 1.00 & $0\%$ \\
Nublado      & 0.95 & $-5\%$ \\
Llovizna     & 0.88 & $-12\%$ \\
Lluvia intensa & 0.75 & $-25\%$ \\
Tormenta     & 0.60 & $-40\%$ \\
\bottomrule
\end{tabular}
\caption{Factor climático multiplicativo sobre parámetros de velocidad.}
\end{table}

% ════════════════════════════════════════════════════════════════════
\section{Marco Teórico: Cadenas de Markov}
% ════════════════════════════════════════════════════════════════════

\subsection{Definición y propiedad de Markov}

\begin{cajadefinicion}{Cadena de Markov de tiempo discreto}
Sea $\{X_t\}_{t \geq 0}$ un proceso estocástico con espacio de estados finito
$\mathcal{S} = \{0, 1, \ldots, k-1\}$. Se dice que es una
\textbf{cadena de Markov de tiempo discreto} si satisface la
\textbf{propiedad de Markov}:
\begin{equation}
    P(X_{t+1} = j \mid X_t = i,\; X_{t-1},\ldots, X_0)
    = P(X_{t+1} = j \mid X_t = i)
    \quad \forall\; i, j \in \mathcal{S},\; t \geq 0
    \label{eq:markov_property}
\end{equation}
\end{cajadefinicion}

La propiedad~\eqref{eq:markov_property} establece que el estado futuro depende
\emph{únicamente} del estado presente. Esta hipótesis de independencia condicional
hace el modelo computacionalmente tratable y es una aproximación razonable para
el tráfico urbano: el estado del flujo vial en $t+1$ está principalmente
determinado por el estado en $t$, no por toda la historia pasada.

\subsection{Espacio de estados en UrbanFlow CDMX}

Definimos $k=3$ estados mutuamente excluyentes:

\begin{table}[h!]
\centering
\renewcommand{\arraystretch}{1.3}
\begin{tabular}{clcc}
\toprule
\textbf{ID} & \textbf{Estado} &
\textbf{Rango de velocidad (ZMVM)} & \textbf{Descripción} \\
\midrule
0 & \textcolor{fluido}{\textbf{Fluido}}        & 20 – 80 km/h & Circulación normal \\
1 & \textcolor{lento}{\textbf{Lento}}          & 5 – 35 km/h  & Reducción perceptible \\
2 & \textcolor{congestionado}{\textbf{Congestionado}} & 2 – 15 km/h & Tráfico detenido \\
\bottomrule
\end{tabular}
\caption{Espacio de estados del tráfico en el modelo de Markov.}
\end{table}

\subsection{Matriz de transición}

La dinámica de la cadena queda completamente descrita por la
\textbf{matriz de transición estocástica} $\mathbf{P} \in \mathbb{R}^{3 \times 3}$:
\begin{equation}
    p_{ij} = P(X_{t+1} = j \mid X_t = i), \qquad i,j \in \{0,1,2\}
    \label{eq:transition_prob}
\end{equation}

\begin{equation}
    \mathbf{P} =
    \begin{pmatrix}
        p_{00} & p_{01} & p_{02} \\
        p_{10} & p_{11} & p_{12} \\
        p_{20} & p_{21} & p_{22}
    \end{pmatrix}
    \label{eq:matrix_P}
\end{equation}

\begin{propiedad}[Estocásticidad]
Cada fila de $\mathbf{P}$ es una distribución de probabilidad:
\begin{equation}
    \sum_{j=0}^{2} p_{ij} = 1 \quad \forall\; i \in \{0,1,2\}
    \label{eq:stochastic}
\end{equation}
\end{propiedad}

\subsection{Distribución en $k$ pasos}

Sea $\bm{\pi}^{(t)} = \bigl(P(X_t=0),\, P(X_t=1),\, P(X_t=2)\bigr)$
el vector de distribución en el instante $t$. La distribución en $t+k$ es:
\begin{equation}
    \bm{\pi}^{(t+k)} = \bm{\pi}^{(t)} \cdot \mathbf{P}^k
    \label{eq:k_step}
\end{equation}

Esto permite predecir, dado el estado actual del tráfico, la distribución
de probabilidad sobre los estados en $k$ minutos (horizonte de predicción).

\subsection{Distribución estacionaria}

\begin{cajadefinicion}{Distribución estacionaria}
Un vector de probabilidad $\bm{\pi}^* = (\pi^*_0, \pi^*_1, \pi^*_2)$ es
la \textbf{distribución estacionaria} de la cadena si satisface:
\begin{equation}
    \bm{\pi}^* \cdot \mathbf{P} = \bm{\pi}^*, \qquad
    \sum_{i=0}^{2} \pi^*_i = 1, \qquad \pi^*_i \geq 0
    \label{eq:stationary}
\end{equation}
\end{cajadefinicion}

Para una cadena \textbf{ergódica} (irreducible y aperiódica),
$\bm{\pi}^*$ es única y representa la fracción de tiempo que el corredor
pasa en cada estado a largo plazo:
\begin{equation}
    \lim_{k \to \infty} \mathbf{P}^k = \mathbf{1} \cdot \bm{\pi}^*
    \quad (\text{convergencia geométrica})
    \label{eq:ergodic_limit}
\end{equation}

En UrbanFlow CDMX, $\bm{\pi}^*$ se calcula encontrando el vector propio
asociado al valor propio $\lambda = 1$ de $\mathbf{P}^\top$:
\begin{equation}
    \mathbf{P}^\top \bm{v} = \bm{v}, \qquad \bm{\pi}^* = \frac{\bm{v}}{\|\bm{v}\|_1}
    \label{eq:eigenvector}
\end{equation}

\subsection{Estimación de la matriz (aprendizaje desde datos)}

Dada una serie histórica de estados $x_0, x_1, \ldots, x_n$, el estimador
de \textbf{máxima verosimilitud} de $p_{ij}$ es:
\begin{equation}
    \hat{p}_{ij} = \frac{n_{ij}}{\displaystyle\sum_{j'} n_{ij'}}
    \label{eq:mle}
\end{equation}
donde $n_{ij} = \#\{t : x_t = i,\; x_{t+1} = j\}$ es el conteo de transiciones.

Para evitar probabilidades cero en estados poco frecuentes se aplica
\textbf{suavizado de Laplace} con constante $\varepsilon = 10^{-6}$:
\begin{equation}
    \hat{p}_{ij} = \frac{n_{ij} + \varepsilon}
                        {\displaystyle\sum_{j'} n_{ij'} + k\,\varepsilon}
    \label{eq:laplace}
\end{equation}

% ════════════════════════════════════════════════════════════════════
\section{Implementación de la Cadena de Markov}
% ════════════════════════════════════════════════════════════════════

\subsection{Clase \texttt{MarkovTrafficChain}}

El módulo \texttt{src/simulation/markov\_chain.py} implementa la clase con
interfaz estilo scikit-learn: \texttt{fit()} $\to$ \texttt{predict\_distribution()}
$\to$ \texttt{simulate()}.

\subsubsection*{Ajuste desde datos históricos}

\begin{lstlisting}[caption={Estimación de la matriz de transición por conteo de transiciones.}]
def fit(self, series: np.ndarray | pd.Series) -> "MarkovTrafficChain":
    arr = _validar_serie(series)

    counts = np.full((N_ESTADOS, N_ESTADOS), self.suavizado)  # epsilon = 1e-6
    for origen, destino in zip(arr[:-1], arr[1:]):
        counts[origen, destino] += 1

    # Normalización fila a fila (ec. MLE con suavizado de Laplace)
    totales = counts.sum(axis=1, keepdims=True)
    filas_cero = (totales == 0).flatten()
    totales[filas_cero] = 1.0
    matriz = counts / totales
    matriz[filas_cero] = 1.0 / N_ESTADOS   # distribución uniforme de respaldo

    self.transition_matrix_ = matriz
    self.n_transitions_ = len(arr) - 1
    return self
\end{lstlisting}

El algoritmo recorre la serie una sola vez ($\mathcal{O}(n)$ en tiempo,
$\mathcal{O}(k^2)$ en espacio con $k=3$ constante).

\subsubsection*{Ejemplo de matriz calibrada para hora pico en la ZMVM}

\begin{equation}
    \hat{\mathbf{P}} =
    \begin{pmatrix}
        0.72 & 0.24 & 0.04 \\
        0.18 & 0.61 & 0.21 \\
        0.05 & 0.28 & 0.67
    \end{pmatrix}
    \begin{matrix}
        \leftarrow \text{desde Fluido} \\
        \leftarrow \text{desde Lento} \\
        \leftarrow \text{desde Congestionado}
    \end{matrix}
    \label{eq:example_matrix}
\end{equation}

\begin{observacion}
Los elementos diagonales $p_{ii}$ son los mayores de cada fila, reflejando la
\textbf{inercia del tráfico}: el estado más probable en el próximo instante es
el mismo que el actual. La probabilidad de mejorar
($\text{Congestionado} \to \text{Fluido}$, $p_{20} = 0.05$) es considerablemente
menor que la de empeorar ($\text{Fluido} \to \text{Congestionado}$, $p_{02} = 0.04$
vs. $p_{20} = 0.05$), consistente con la asimetría del tráfico observada en la ZMVM.
\end{observacion}

\subsubsection*{Predicción de distribución en $k$ pasos}

\begin{lstlisting}[caption={Implementación de la ecuación~\eqref{eq:k_step}.}]
def predict_distribution(self, estado_inicial: int, pasos: int) -> np.ndarray:
    v = np.zeros(N_ESTADOS)
    v[estado_inicial] = 1.0                   # vector one-hot del estado inicial
    P_k = np.linalg.matrix_power(self.transition_matrix_, pasos)
    return v @ P_k                            # distribucion en t + pasos
\end{lstlisting}

\subsubsection*{Distribución estacionaria}

\begin{lstlisting}[caption={Cálculo de $\bm{\pi}^*$ por descomposición espectral, ecuación~\eqref{eq:eigenvector}.}]
def steady_state(self) -> np.ndarray:
    valores, vectores = np.linalg.eig(self.transition_matrix_.T)
    idx = np.argmin(np.abs(valores - 1.0))   # valor propio lambda = 1
    pi = np.real(vectores[:, idx])
    pi = np.abs(pi)
    return pi / pi.sum()                      # normalizar a distribucion de prob.
\end{lstlisting}

% ════════════════════════════════════════════════════════════════════
\section{Marco Teórico: Simulación Monte Carlo}
% ════════════════════════════════════════════════════════════════════

\subsection{Fundamento matemático}

La \textbf{simulación Monte Carlo} es un método de integración numérica estocástica
que estima distribuciones de probabilidad de variables de salida mediante muestreo
repetido. Fue formalizado por Stanisław Ulam y John von Neumann (Proyecto Manhattan,
1947) y toma su nombre del Casino de Montecarlo.

\begin{cajadefinicion}{Estimador Monte Carlo}
Sea $Y = g(\mathbf{X})$ una variable de salida función de entradas aleatorias.
El estimador Monte Carlo de $E[Y]$ con $N$ muestras independientes es:
\begin{equation}
    \hat{\mu}_N = \frac{1}{N} \sum_{i=1}^{N} g(\mathbf{X}_i),
    \qquad
    \hat{\mu}_N \xrightarrow[N\to\infty]{\text{c.s.}} E[Y]
    \quad (\text{Ley de los Grandes Números})
    \label{eq:mc_estimator}
\end{equation}
El error estándar del estimador decrece como:
\begin{equation}
    \text{SE}[\hat{\mu}_N] = \frac{\sigma_Y}{\sqrt{N}}
    \label{eq:mc_error}
\end{equation}
\end{cajadefinicion}

Con $N = 10{,}000$ simulaciones, el error estándar es
$\sigma_Y / 100$, garantizando estimaciones de percentiles con
precisión de $\pm 0.5$ min para los rangos de tiempo típicos de la ZMVM.

\subsection{Distribución de velocidades por estado}

Para cada estado de tráfico $s \in \{0, 1, 2\}$, la velocidad instantánea
se modela como una \textbf{Normal truncada}:
\begin{equation}
    V \mid s \;\sim\; \mathcal{TN}\!\left(\mu_s,\, \sigma_s^2;\; [v^\text{min}_s,\, v^\text{max}_s]\right)
    \label{eq:truncated_normal}
\end{equation}

Los parámetros están calibrados con datos del TomTom Traffic Index CDMX
y aforos SEMOVI:

\begin{table}[h!]
\centering
\renewcommand{\arraystretch}{1.3}
\begin{tabular}{lcccc}
\toprule
\textbf{Estado} & $\mu_s$ (km/h) & $\sigma_s$ (km/h) & $v^\text{min}_s$ & $v^\text{max}_s$ \\
\midrule
\textcolor{fluido}{\textbf{Fluido}}               & 40.0 & 8.0 & 20.0 & 80.0 \\
\textcolor{lento}{\textbf{Lento}}                 & 18.0 & 5.0 &  5.0 & 35.0 \\
\textcolor{congestionado}{\textbf{Congestionado}} &  7.0 & 3.0 &  2.0 & 15.0 \\
\bottomrule
\end{tabular}
\caption{Parámetros de la distribución Normal truncada por estado de tráfico.}
\end{table}

\subsection{Modelo de tiempo de viaje}

Sea $d$ la distancia del recorrido en km y $\Delta t$ el paso temporal
(1 minuto). Para la trayectoria $i$-ésima, la simulación calcula:

\begin{align}
    s^{(i)}_t &\sim \text{Markov}\!\left(s^{(i)}_{t-1},\; \mathbf{P}\right)
        && \text{(estado en el paso } t \text{)} \label{eq:mc_state} \\
    v^{(i)}_t &\sim \mathcal{TN}\!\left(\mu_{s^{(i)}_t},\; \sigma_{s^{(i)}_t};\;
        [v^\text{min}, v^\text{max}]\right)
        && \text{(velocidad en km/h)} \label{eq:mc_velocity} \\
    \delta d^{(i)}_t &= v^{(i)}_t \cdot \frac{\Delta t}{60}
        && \text{(distancia cubierta en km)} \label{eq:mc_dist_step}
\end{align}

El tiempo de llegada de la trayectoria $i$ es:
\begin{equation}
    T^{(i)} = \min\!\left\{t \;\Big|\; \sum_{\tau=1}^{t} \delta d^{(i)}_\tau \geq d\right\}
    \label{eq:mc_arrival}
\end{equation}

Repitiendo \eqref{eq:mc_state}–\eqref{eq:mc_arrival} para $i = 1, \ldots, N$
se obtiene la distribución empírica $\{T^{(1)}, \ldots, T^{(N)}\}$
de la cual se extraen los percentiles:
\begin{equation}
    P_{q} = Q_q\!\left(\{T^{(i)}\}_{i=1}^{N}\right), \quad q \in \{10, 50, 90\}
    \label{eq:percentiles}
\end{equation}

% ════════════════════════════════════════════════════════════════════
\section{Implementación del Motor Monte Carlo}
% ════════════════════════════════════════════════════════════════════

\subsection{Clase \texttt{MonteCarloEngine}}

El módulo \texttt{src/simulation/monte\_carlo.py} implementa el motor con
simulación completamente \textbf{vectorizada} sobre el eje de trayectorias
($N \times T$ operaciones simultáneas via NumPy), eliminando cualquier
bucle de Python sobre las simulaciones.

\begin{algorithm}[H]
\SetAlgoLined
\KwIn{Cadena ajustada $\hat{\mathbf{P}}$, distancia $d$ km, estado inicial $s_0$,
      $N = 10{,}000$, $T_\text{max} = 480$ pasos}
\KwOut{$P_{10}$, $P_{50}$, $P_{90}$, $\bar{T}$, $\sigma_T$}
\BlankLine
Precomputar $\mathbf{P}_\text{cum} \leftarrow \texttt{cumsum}(\hat{\mathbf{P}},\; \text{axis}=1)$\;
\BlankLine
\tcp{Paso 1: Generar trayectorias de estados (vectorizado)}
$\mathbf{S} \leftarrow \mathbf{0}_{N \times T_\text{max}}$ \tcp*{matriz de estados int8}
$\mathbf{S}[:, 0] \leftarrow s_0$\;
\For{$t = 1$ \KwTo $T_\text{max} - 1$}{
    $\mathbf{u} \sim \mathcal{U}[0,1)^N$ \tcp*{N números uniformes simultáneos}
    $\mathbf{S}[:,t] \leftarrow \bigl(\mathbf{u}_{:,\text{new}} \geq \mathbf{P}_\text{cum}[\mathbf{S}[:,t-1]]\bigr)\text{.sum}(\text{axis}=1)$\;
}
\BlankLine
\tcp{Paso 2: Muestrear velocidades por estado (máscara booleana)}
$\mathbf{V} \leftarrow \mathbf{0}_{N \times T_\text{max}}$\;
\For{$s \in \{0, 1, 2\}$}{
    $\text{mask} \leftarrow (\mathbf{S} = s)$\;
    $\mathbf{V}[\text{mask}] \sim \mathcal{TN}(\mu_s, \sigma_s; [v^\text{min}_s, v^\text{max}_s])$\;
}
\BlankLine
\tcp{Paso 3: Distancia acumulada e interpolación sub-paso}
$\mathbf{D} \leftarrow \texttt{cumsum}(\mathbf{V} \cdot \frac{\Delta t}{60},\; \text{axis}=1)$\;
$\text{idx} \leftarrow \texttt{argmax}(\mathbf{D} \geq d,\; \text{axis}=1)$\;
$\mathbf{T} \leftarrow \bigl(\text{idx} + \text{fracción lineal}\bigr) \cdot \Delta t$\;
\BlankLine
\Return $P_{10}, P_{50}, P_{90} \leftarrow \texttt{percentile}(\mathbf{T},\; [10, 50, 90])$\;
\caption{Motor Monte Carlo vectorizado de UrbanFlow CDMX}
\end{algorithm}

\subsection{Generación vectorizada de estados de Markov}

El paso más costoso es la generación de $N \times T_\text{max}$ estados.
En lugar de llamar a \texttt{rng.choice()} $N$ veces en bucle, se usa el
\textbf{método de la transformada inversa vectorizado}:

\begin{lstlisting}[caption={Generación de $N$ estados simultáneos por muestreo de CDF acumulada.}]
def _simular_estados(self, estado_inicial: int) -> np.ndarray:
    n, T = self.n_simulaciones, self.max_pasos
    estados = np.empty((n, T), dtype=np.int8)
    estados[:, 0] = estado_inicial

    for t in range(1, T):
        actual = estados[:, t - 1]              # (N,) estado actual por simulacion
        u = self._rng.random(n)                 # (N,) numeros uniformes [0, 1)

        # P_cumsum[actual] -> (N, 3): CDF acumulada por simulacion
        # Contar cuantos umbrales son < u determina el siguiente estado
        siguiente = (u[:, np.newaxis] >= self._P_cumsum[actual]).sum(axis=1)
        estados[:, t] = np.clip(siguiente, 0, N_ESTADOS - 1).astype(np.int8)

    return estados  # shape: (N, T_max)
\end{lstlisting}

\subsubsection*{Ejemplo del método de transformada inversa}

Para el estado \textbf{Lento} ($i=1$), la CDF acumulada de la fila 1 de
$\hat{\mathbf{P}}$ es:
\[
\mathbf{P}_\text{cum}[1] = [0.18,\; 0.79,\; 1.00]
\]

Dado $u \sim \mathcal{U}[0,1)$, el siguiente estado se asigna según:
\[
j^* = \sum_{j=0}^{2} \mathbf{1}\!\left[u \geq \mathbf{P}_\text{cum}[1,j]\right]
\]

\begin{table}[h!]
\centering
\renewcommand{\arraystretch}{1.3}
\begin{tabular}{cccc}
\toprule
$u$ & $u \geq 0.18$ & $u \geq 0.79$ & $j^* = $ Estado siguiente \\
\midrule
0.10 & \xmark & \xmark & $0$ (Fluido) \\
0.45 & \checkmark & \xmark & $1$ (Lento) \\
0.85 & \checkmark & \checkmark & $2$ (Congestionado) \\
\bottomrule
\end{tabular}
\caption{Método de transformada inversa para muestreo del estado siguiente.
         \checkmark\,= condición cumplida;\; \xmark\,= condición no cumplida.}
\end{table}

\subsection{Cálculo del tiempo de llegada con interpolación lineal}

Para evitar el redondeo al minuto más cercano, se aplica interpolación lineal
sub-paso. Sea $t^* = \texttt{argmax}_{t}(D^{(i)}_t \geq d)$ el primer paso
donde se supera la distancia objetivo:
\begin{equation}
    T^{(i)} = \left(t^* - 1 + \frac{d - D^{(i)}_{t^*-1}}{D^{(i)}_{t^*} - D^{(i)}_{t^*-1}}\right)
              \cdot \Delta t \quad \text{(minutos)}
    \label{eq:interpolation}
\end{equation}

Esta expresión da precisión sub-minuto sin necesidad de reducir $\Delta t$.

\subsection{Estructura del resultado}

\begin{lstlisting}[caption={Dataclass \texttt{ResultadoSimulacion}: salida completa del motor.}]
@dataclass
class ResultadoSimulacion:
    tiempos_minutos: np.ndarray  # distribucion completa (N = 10,000 valores)
    p10:  float                  # percentil 10: viaje rapido
    p50:  float                  # percentil 50: mediana (mas probable)
    p90:  float                  # percentil 90: viaje pesimista
    media: float                 # tiempo esperado E[T]
    std:   float                 # desviacion estandar sigma_T
    banda_incertidumbre: float   # P90 - P10 (amplitud de incertidumbre)
    n_recortadas: int            # trayectorias que excedieron max_pasos
\end{lstlisting}

\begin{cajaresultado}{Ejemplo de resultado real}
Para un recorrido de 12.5 km con estado inicial \textbf{Lento} y
lluvia leve ($\phi = 0.88$):
\[
P_{10} = 28.3\;\text{min}, \quad
P_{50} = 37.1\;\text{min}, \quad
P_{90} = 52.8\;\text{min}, \quad
\text{Banda} = 24.5\;\text{min}
\]
\end{cajaresultado}

% ════════════════════════════════════════════════════════════════════
\section{Pipeline Integrador en Tiempo Real}
\label{sec:pipeline}
% ════════════════════════════════════════════════════════════════════

\subsection{Inferencia del estado de tráfico}

El estado inicial $s_0$ de la cadena de Markov se infiere del ratio de
congestión promedio del corredor $\bar{r}$ (ecuación~\eqref{eq:ratio}):
\begin{equation}
    s_0 =
    \begin{cases}
        0 \;(\text{Fluido})        & \text{si } \bar{r} \geq 0.75 \\
        1 \;(\text{Lento})         & \text{si } 0.45 \leq \bar{r} < 0.75 \\
        2 \;(\text{Congestionado}) & \text{si } \bar{r} < 0.45
    \end{cases}
    \label{eq:state_inference}
\end{equation}

Los umbrales $\{0.75, 0.45\}$ están calibrados empíricamente para la ZMVM
y son consistentes con la clasificación del TomTom Traffic Index.

\subsection{Ajuste climático de los parámetros de velocidad}

El factor climático $\phi$ modifica las distribuciones de velocidad antes de
iniciar la simulación Monte Carlo:
\begin{equation}
    \mu'_s = \phi \cdot \mu_s, \quad
    v'^{\text{min}}_s = \phi \cdot v^{\text{min}}_s, \quad
    v'^{\text{max}}_s = \phi \cdot v^{\text{max}}_s
    \label{eq:climate_adjustment}
\end{equation}
La desviación estándar $\sigma_s$ no se modifica: la dispersión del tráfico
no cambia por la lluvia, sólo disminuye la velocidad media.

\subsection{Flujo completo de predicción}

\begin{enumerate}
    \item \textbf{TomTom Routing API} $\to$ $d$ km, waypoints (geometría del mapa)
    \item \textbf{TomTom Traffic Stats} $\to$ $\{r_k\}_{k=1}^{K}$, $\bar{r}$ promedio
    \item Inferir $s_0$ con \eqref{eq:state_inference}
    \item \textbf{OWM API} $\to$ condición climática $\to$ factor $\phi$
    \item Ajustar parámetros de velocidad con \eqref{eq:climate_adjustment}
    \item \texttt{ConsultaViaje}$(d, s_0)$ $\to$ \texttt{MonteCarloEngine.correr()}
    \item $\Rightarrow$ $P_{10}$, $P_{50}$, $P_{90}$ en $\approx 500\;\text{ms}$
\end{enumerate}

% ════════════════════════════════════════════════════════════════════
\section{Resultados y Validación}
% ════════════════════════════════════════════════════════════════════

\subsection{Convergencia del estimador Monte Carlo}

A partir del error estándar del estimador \eqref{eq:mc_error} y de los rangos
de tiempo típicos de la ZMVM ($\sigma_T \approx 15$–$20$ min):

\begin{table}[h!]
\centering
\renewcommand{\arraystretch}{1.3}
\begin{tabular}{rcc}
\toprule
$N$ & Error est. $P_{50}$ (min) & Error est. $P_{90}$ (min) \\
\midrule
100       & $\pm 5.2$ & $\pm 8.1$ \\
1\,000    & $\pm 1.6$ & $\pm 2.6$ \\
10\,000   & $\pm 0.5$ & $\pm 0.8$ \\
100\,000  & $\pm 0.2$ & $\pm 0.3$ \\
\bottomrule
\end{tabular}
\caption{Error de estimación de percentiles en función de $N$. Balance óptimo
         con $N = 10{,}000$: precisión $\pm 0.5$ min, cómputo $< 1$ s.}
\end{table}

\subsection{Bandas de incertidumbre por estado de tráfico}

Para un corredor de 15 km bajo distintas condiciones iniciales:

\begin{table}[h!]
\centering
\renewcommand{\arraystretch}{1.3}
\begin{tabular}{lcccc}
\toprule
\textbf{Estado inicial} & $P_{10}$ (min) & $P_{50}$ (min) & $P_{90}$ (min) & Banda (min) \\
\midrule
\textcolor{fluido}{\textbf{Fluido}}               & 22 & 28 & 38 & 16 \\
\textcolor{lento}{\textbf{Lento}}                 & 30 & 42 & 61 & 31 \\
\textcolor{congestionado}{\textbf{Congestionado}} & 48 & 72 & 105 & 57 \\
\bottomrule
\end{tabular}
\caption{Bandas de incertidumbre $P_{10}/P_{50}/P_{90}$ para un recorrido de 15 km.}
\end{table}

La banda $P_{90} - P_{10}$ se amplía con la congestión, capturando la mayor
variabilidad del tráfico denso: cuando el tráfico está fluido la banda es de
16 min, pero en congestión alcanza 57 min.

\subsection{Métrica objetivo MAPE}

La métrica de éxito del proyecto es:
\begin{equation}
    \text{MAPE} = \frac{1}{n} \sum_{k=1}^{n}
                  \frac{|T^\text{real}_k - T^{P_{50}}_k|}{T^\text{real}_k}
                  \times 100\% < 15\%
    \label{eq:mape}
\end{equation}
donde $T^{P_{50}}_k$ es la mediana de la distribución Monte Carlo para la
consulta $k$. El umbral del 15\% fue establecido como objetivo del diplomado
para horizontes de predicción de 30 minutos.

% ════════════════════════════════════════════════════════════════════
\section{Conclusiones y Trabajo Futuro}
% ════════════════════════════════════════════════════════════════════

\subsection{Conclusiones}

\begin{enumerate}
    \item \textbf{La modelación estocástica supera al enfoque determinístico.}
    Las bandas de incertidumbre $P_{10}/P_{90}$ proveen información accionable que
    un tiempo puntual no puede ofrecer: el usuario puede planearse con el nivel
    de certeza que necesite.

    \item \textbf{Las cadenas de Markov son un modelo adecuado para la dinámica
    del tráfico.} La propiedad de Markov~\eqref{eq:markov_property} es una
    aproximación razonable: el estado vial en $t+1$ está correlacionado principalmente
    con el estado en $t$. Los elementos diagonales dominantes en $\hat{\mathbf{P}}$
    confirman la inercia del tráfico observada en la ZMVM.

    \item \textbf{Monte Carlo con $N = 10{,}000$ trayectorias es suficiente.}
    El error estándar del estimador de $P_{50}$ es $\pm 0.5$ min, con tiempo
    de cómputo menor a 1 segundo en CPU sin GPU.

    \item \textbf{El factor climático mejora la calibración.} El ajuste
    multiplicativo~\eqref{eq:climate_adjustment} captura el efecto de la lluvia
    sobre las distribuciones de velocidad de forma simple y sin sobreajuste.

    \item \textbf{La TomTom Routing API es una mejora sustancial.} La red vial
    urbana tiene tortuosidad media del 40\% (factor $\kappa = 1.4$); usar la
    distancia real por carretera reduce el sesgo de la predicción.
\end{enumerate}

\subsection{Trabajo futuro}

\begin{table}[h!]
\centering
\renewcommand{\arraystretch}{1.3}
\begin{tabularx}{\textwidth}{llX}
\toprule
\textbf{Prioridad} & \textbf{Mejora} & \textbf{Impacto esperado} \\
\midrule
Alta & Matrices de Markov diferenciadas por hora y corredor &
Capturar patrón hora-pico vs. hora-valle \\
Alta & Pipeline de entrenamiento con datos C5 CDMX (+1M registros) &
Estimaciones $\hat{p}_{ij}$ estadísticamente robustas \\
Media & API REST FastAPI + caché Redis &
Reducir latencia a $< 100$ ms \\
Media & Integración con Metrobús GTFS-RT &
Predicción multimodal (auto + transporte público) \\
Baja & Motor Monte Carlo en GPU (CuPy) &
Reducir tiempo de simulación a $< 50$ ms \\
Baja & Corrector de sesgo XGBoost post-Monte Carlo &
Reducir MAPE residual \\
\bottomrule
\end{tabularx}
\caption{Líneas de trabajo futuro del proyecto UrbanFlow CDMX.}
\end{table}

% ════════════════════════════════════════════════════════════════════
\section*{Referencias}
% ════════════════════════════════════════════════════════════════════
\addcontentsline{toc}{section}{Referencias}

\begin{enumerate}
    \item Norris, J. R. (1997). \textit{Markov Chains}. Cambridge University Press.

    \item Rubinstein, R. Y., \& Kroese, D. P. (2016).
    \textit{Simulation and the Monte Carlo Method} (3rd ed.). Wiley.

    \item TomTom Traffic Index 2024 — Ciudad de México.
    \texttt{https://www.tomtom.com/traffic-index/}

    \item SEMOVI (2023). \textit{Aforos vehiculares en la Zona Metropolitana del
    Valle de México}. Secretaría de Movilidad, Gobierno de la Ciudad de México.

    \item Harris, C. R. et al. (2020). Array programming with NumPy.
    \textit{Nature}, 585, 357–362.
    \texttt{https://doi.org/10.1038/s41586-020-2649-2}

    \item Van Lint, J. W. C., \& Hoogendoorn, S. P. (2010).
    A Robust and Efficient Method for Fusing Heterogeneous Data from Traffic Sensors.
    \textit{Computer-Aided Civil and Infrastructure Engineering}, 25(8), 596–612.

    \item Sinnott, R. W. (1984). Virtues of the Haversine.
    \textit{Sky and Telescope}, 68(2), 159.

    \item Streamlit Documentation (2024).
    \textit{Streamlit — The fastest way to build data apps.}
    \texttt{https://docs.streamlit.io}
\end{enumerate}

% ════════════════════════════════════════════════════════════════════
\appendix
% ════════════════════════════════════════════════════════════════════

\section{Parámetros de Configuración del Motor}

\begin{table}[h!]
\centering
\renewcommand{\arraystretch}{1.3}
\begin{tabular}{lll}
\toprule
\textbf{Parámetro} & \textbf{Valor por defecto} & \textbf{Descripción} \\
\midrule
\texttt{n\_simulaciones}    & 10{,}000       & Número de trayectorias Monte Carlo \\
\texttt{paso\_minutos}      & 1.0            & Resolución temporal $\Delta t$ (min) \\
\texttt{max\_pasos}         & 480            & Horizonte máximo (8 horas) \\
\texttt{suavizado}          & $10^{-6}$      & Constante de Laplace $\varepsilon$ \\
\texttt{FACTOR\_TORTUOSIDAD} & 1.4           & Factor $\kappa$ Haversine $\to$ carretera \\
\texttt{UMBRAL\_FLUIDO}     & 0.75           & Umbral de ratio para estado Fluido \\
\texttt{UMBRAL\_LENTO}      & 0.45           & Umbral de ratio para estado Lento \\
\bottomrule
\end{tabular}
\caption{Parámetros configurables del motor estocástico.}
\end{table}

\section{Glosario}

\begin{description}
    \item[Cadena de Markov] Proceso estocástico donde el estado futuro depende
        únicamente del estado presente (propiedad de Markov).
    \item[Distribución estacionaria ($\bm{\pi}^*$)] Vector de probabilidades
        al que converge la cadena de Markov en el largo plazo.
    \item[Monte Carlo] Método de estimación basado en muestreo aleatorio masivo.
    \item[Normal truncada] Distribución normal restringida a un intervalo
        $[v^\text{min}, v^\text{max}]$.
    \item[$P_{10} / P_{50} / P_{90}$] Percentiles 10, 50 y 90 de la distribución
        simulada de tiempos de viaje.
    \item[Ratio de congestión ($r$)] Cociente $v_\text{actual} / v_\text{libre}$;
        mide el grado de congestión del segmento vial.
    \item[MAPE] \textit{Mean Absolute Percentage Error}; error porcentual
        absoluto medio.
    \item[ZMVM] Zona Metropolitana del Valle de México.
    \item[WGS84] Sistema de coordenadas geográficas estándar (GPS, EPSG:4326).
    \item[EPSG:6372] Sistema de referencia cartográfico oficial para México
        (para cálculos de distancia en metros).
\end{description}

\end{document}
